% GENERATED FILE. Edit canonical lessons, then run npm run generate:books.
% canonical-source-hash: fnv1a64:6582b4ba25900616
% canonical-lessons: RU-C39-myt, RU-C39-gotovit, RU-C39-pokupat, RU-C39-prodavat, RU-C39-plavat

\chapter{Wash, Cook, Buy, Sell, Swim}
\label{ch:ru-wash-cook-buy-sell-swim}
\hlchaptermodality{39}

\begin{chapteropening}
\textbf{By the end of this chapter:} \emph{I can say to wash, to cook, to buy, to sell and to swim.}

Complete the last lesson of chapter 39: I can say to wash, to cook, to buy, to sell and to swim.
\end{chapteropening}

\section[\texorpdfstring{\ru{мыть} (myt)}{мыть (myt)}]{\ru{мыть} (myt) — to wash}

\label{lesson:RU-C39-myt}



\begin{quote}
Before the new one: say the Russian for to laugh, then the Russian for to cry.
\end{quote}

\subsection*{You'll want to know: \ru{мыть}}
\textbf{\ru{мыть}} (\emph{myt}) — \textquotedblleft{}to wash\textquotedblright{}.

\begin{cousinweb}
\textbf{\ru{мыть}} $\leftarrow$ Proto-Slavic \emph{*myti}, present \emph{*myjǫ} — inherited, a cousin of Lithuanian \emph{máudyti}, \textquotedblleft{}to bathe, to wash\textquotedblright{}, and Latvian \emph{maût}, \textquotedblleft{}to swim\textquotedblright{} (Vasmer).
\end{cousinweb}

\begin{grammarlens}[title={What you've built}]
One more word for this chapter.
\end{grammarlens}

\subsection*{Your turn}
\begin{itemize}
\raggedright
  \item \emph{Say it:} \emph{myt}
  \item \emph{Say it:} \emph{myt}, once more
  \item \emph{Say it:} say \emph{myt}
  \item \emph{Recall:} say the Russian for to laugh, then the Russian for to cry, then say \emph{myt} again
\end{itemize}

\begin{tcolorbox}[breakable,colback=teal!4,colframe=teal!35!black,title={Before you move on}]
What does \ru{мыть} mean? (\textquotedblleft{}To wash\textquotedblright{}.) Say it once more.
\end{tcolorbox}
\section[\texorpdfstring{\ru{готовить} (gotóvit)}{готовить (gotóvit)}]{\ru{готовить} (gotóvit) — to cook, to prepare}

\label{lesson:RU-C39-gotovit}



\begin{quote}
Before the new one: say the Russian for to cry, then the Russian for to wash.
\end{quote}

\subsection*{You'll want to know: \ru{готовить}}
\textbf{\ru{готовить}} (\emph{gotóvit}) — \textquotedblleft{}to cook, to prepare\textquotedblright{}.

\begin{cousinweb}
\textbf{\ru{готовить}} is \emph{\ru{готовый}}, \textquotedblleft{}ready\textquotedblright{}, made into a verb: \textquotedblleft{}to make ready\textquotedblright{}, and so \textquotedblleft{}to cook\textquotedblright{}. \emph{\ru{готовый}} goes back to Proto-Slavic \emph{*gotovŭ}; Vasmer relates it to Albanian \emph{gatuanj}, \textquotedblleft{}I cook\textquotedblright{}, and rejects a loan from Gothic.
\end{cousinweb}

\begin{grammarlens}[title={What you've built}]
One more word for this chapter.
\end{grammarlens}

\subsection*{Your turn}
\begin{itemize}
\raggedright
  \item \emph{Say it:} \emph{gotóvit}
  \item \emph{Say it:} \emph{gotóvit}, once more
  \item \emph{Say it:} say \emph{gotóvit}
  \item \emph{Recall:} say the Russian for to cry, then the Russian for to wash, then say \emph{gotóvit} again
\end{itemize}

\begin{tcolorbox}[breakable,colback=teal!4,colframe=teal!35!black,title={Before you move on}]
What does \ru{готовить} mean? (\textquotedblleft{}To cook, to prepare\textquotedblright{}.) Say it once more.
\end{tcolorbox}
\section[\texorpdfstring{\ru{покупать} (pokupát)}{покупать (pokupát)}]{\ru{покупать} (pokupát) — to buy}

\label{lesson:RU-C39-pokupat}



\begin{quote}
Before the new one: say the Russian for to wash, then the Russian for to cook.
\end{quote}

\subsection*{You'll want to know: \ru{покупать}}
\textbf{\ru{покупать}} (\emph{pokupát}) — \textquotedblleft{}to buy\textquotedblright{}.

\begin{cousinweb}
\textbf{\ru{покупать}} is the ongoing partner of \emph{\ru{купить}}, \textquotedblleft{}to buy\textquotedblright{}, with \emph{\ru{по}-}. \emph{\ru{купить}} was borrowed from Gothic \emph{kaupōn}, \textquotedblleft{}to trade\textquotedblright{}, and the Germanic word comes from Latin \emph{caupō}, \textquotedblleft{}innkeeper, petty trader\textquotedblright{} (Vasmer). German \emph{kaufen}, \textquotedblleft{}to buy\textquotedblright{}, has the same source.
\end{cousinweb}

\begin{grammarlens}[title={What you've built}]
One more word for this chapter.
\end{grammarlens}

\subsection*{Your turn}
\begin{itemize}
\raggedright
  \item \emph{Say it:} \emph{pokupát}
  \item \emph{Say it:} \emph{pokupát}, once more
  \item \emph{Say it:} say \emph{pokupát}
  \item \emph{Recall:} say the Russian for to wash, then the Russian for to cook, then say \emph{pokupát} again
\end{itemize}

\begin{tcolorbox}[breakable,colback=teal!4,colframe=teal!35!black,title={Before you move on}]
What does \ru{покупать} mean? (\textquotedblleft{}To buy\textquotedblright{}.) Say it once more.
\end{tcolorbox}
\section[\texorpdfstring{\ru{продавать} (prodavát)}{продавать (prodavát)}]{\ru{продавать} (prodavát) — to sell}

\label{lesson:RU-C39-prodavat}



\begin{quote}
Before the new one: say the Russian for to cook, then the Russian for to buy.
\end{quote}

\subsection*{You'll want to know: \ru{продавать}}
\textbf{\ru{продавать}} (\emph{prodavát}) — \textquotedblleft{}to sell\textquotedblright{}.

\begin{cousinweb}
\textbf{\ru{продавать}} is \emph{\ru{про}-}, \textquotedblleft{}through, forth\textquotedblright{}, with \emph{\ru{давать}}, \textquotedblleft{}to give\textquotedblright{}, on the root of \emph{\ru{дать}}, \textquotedblleft{}to give\textquotedblright{}, which Vasmer relates to Lithuanian \emph{dúoti} and Greek \emph{didōmi}, \textquotedblleft{}to give\textquotedblright{}; Latin \emph{dare} is a cousin too.
\end{cousinweb}

\begin{grammarlens}[title={What you've built}]
One more word for this chapter.
\end{grammarlens}

\subsection*{Your turn}
\begin{itemize}
\raggedright
  \item \emph{Say it:} \emph{prodavát}
  \item \emph{Say it:} \emph{prodavát}, once more
  \item \emph{Say it:} say \emph{prodavát}
  \item \emph{Recall:} say the Russian for to cook, then the Russian for to buy, then say \emph{prodavát} again
\end{itemize}

\begin{tcolorbox}[breakable,colback=teal!4,colframe=teal!35!black,title={Before you move on}]
What does \ru{продавать} mean? (\textquotedblleft{}To sell\textquotedblright{}.) Say it once more.
\end{tcolorbox}
\section[\texorpdfstring{\ru{плавать} (plávat)}{плавать (plávat)}]{\ru{плавать} (plávat) — to swim}

\label{lesson:RU-C39-plavat}



\begin{quote}
Before the new one: say the Russian for to buy, then the Russian for to sell.
\end{quote}

\subsection*{You'll want to know: \ru{плавать}}
\textbf{\ru{плавать}} (\emph{plávat}) — \textquotedblleft{}to swim\textquotedblright{}.

\begin{cousinweb}
\textbf{\ru{плавать}} $\leftarrow$ Proto-Slavic \emph{*plavati} — inherited, a repeated-action form on the root of \emph{\ru{плыть}}, \textquotedblleft{}to swim, to sail\textquotedblright{} (Vasmer). Vasmer relates that root to Sanskrit \emph{plávate}, \textquotedblleft{}floats\textquotedblright{}, Greek \emph{pléō}, \textquotedblleft{}I sail\textquotedblright{}, and Latin \emph{pluit}, \textquotedblleft{}it rains\textquotedblright{}.
\end{cousinweb}

\begin{grammarlens}[title={What you've built}]
One more word for this chapter.
\end{grammarlens}

\subsection*{Your turn}
\begin{itemize}
\raggedright
  \item \emph{Say it:} \emph{plávat}
  \item \emph{Say it:} \emph{plávat}, once more
  \item \emph{Say it:} say \emph{plávat}
  \item \emph{Recall:} say the Russian for to buy, then the Russian for to sell, then say \emph{plávat} again
\end{itemize}

\begin{tcolorbox}[breakable,colback=teal!4,colframe=teal!35!black,title={Before you move on}]
What does \ru{плавать} mean? (\textquotedblleft{}To swim\textquotedblright{}.) Say it once more.
\end{tcolorbox}
